\chapter{Learned look-ahead and the priors that override it}
\label{ch:lookahead}

\begin{objectives}
\begin{itemize}
\item State what "learned look-ahead" claims, and what evidence has been offered for it.
\item Separate what is published, what is architecture-specific, and what this project hopes
  is true.
\item Examine a real position in which all three of Leela's output heads miss a mate in two.
\item Design the suppressed-win experiment properly: probe, controls, and the decisive test.
\item Compute why base rates make a naive suppressed-win detector produce mostly false alarms
  --- and what to do about it.
\end{itemize}
\end{objectives}

\section{The claim}

Everything in Part~\textsc{ii} assumed a clean division of labour: the network has intuition,
the search does the calculating. Learned look-ahead is the claim that this division is not
clean --- that \textbf{inside a single forward pass, with no search at all, the network's
middle layers construct representations of positions several plies in the future, and use
them}.

If true it is a significant thing to know, and it is directly on this project's north star:
it would mean the network's "thinking" contains more structure than its outputs reveal, and
that structure is what a coaching tool would most want to translate.

\begin{researchfinding}
\textbf{Evidence of Learned Look-Ahead in a Chess-Playing Neural Network} (Jenner et al.,
2024) studied a Leela chess network and reported, in summary:
\begin{enumerate}
\item \textbf{Future-state representation.} Activations in middle layers contain information
  about board states several plies ahead; a linear probe recovered aspects of the optimal move
  two turns ahead with accuracy reported up to around $92\%$.
\item \textbf{Look-ahead attention heads.} Specific heads move information along the squares
  of a future tactical sequence, rather than the current one.
\item \textbf{Causal relevance.} Intervening on those activations changes the network's
  immediate move choice --- so the representations are used, not decorative
  (Chapter~\ref{ch:probes}, Trap 3, addressed).
\end{enumerate}
Point 3 is what makes this a strong result rather than a suggestive one.
\end{researchfinding}

\begin{pitfall}
Three limits on how far that carries.

\textbf{(a) It is one network.} The result was obtained on a particular Leela net. Your
\texttt{BT3-768x15x24h} is a different, larger architecture. Interpretability findings are
routinely architecture-specific; the finding is a reason to look, not a licence to assume.

\textbf{(b) "Look-ahead" is a claim about representation, not about search.} The network is
not running a tree. It has learned patterns that encode where pieces will be in common
tactical shapes. That is impressive; it is not the same as calculating, and a tool that says
"the network calculated this line" would be overclaiming.

\textbf{(c) There is a citation to verify.} Your project's research note also lists a paper
titled \emph{Learned Priors Override Look-Ahead in a Chess-Playing Neural Network}. I have not
been able to verify that this is a published paper with that title, and this book will not
treat it as one. The \emph{phenomenon} it names is separately plausible and partly observable
in our own data (§15.2) --- but before that title is cited anywhere in this project's
documents, someone should confirm it exists. Given this project's history with fabricated
sources, treat an unverified citation as a bug.
\end{pitfall}

\section{A real position where every head is wrong}

We do not have to speculate about whether Leela's outputs can miss a decisive tactic. We
measured it in Chapter~\ref{ch:puct}, and it is worth assembling all the evidence in one
place, because this position is the natural first test case for everything that follows.

\begin{center}
\chessboard[setfen=4kb1r/p2n1ppp/4q3/4p1B1/4P3/1Q6/PPP2PPP/2KR4 w k - 0 16,
            boardfontsize=18pt]
\end{center}

\noindent Morphy--Duke of Brunswick \& Count Isouard, Paris 1858, before 16.Qb8+. White to
move; 46 legal moves. Stockfish at depth 30: \textbf{16.Qb8+ Nxb8 17.Rd8\# --- mate in two.}

\begin{realdata}[title={Real engine output: what the network says, with no search}]
\begin{center}
\begin{tabular}{@{}llr@{}}
\toprule
Head & Output & Reading \\
\midrule
Policy & $P(\text{Qb8+}) = 1.60\%$ & the mating move is not among the top few \\
Policy & $P(\text{Qb7}) = 31.93\%$ & prefers a strong but non-forcing move \\
Value  & $V = +0.62$, $d = 0.193$ & "clearly better", not "winning by force" \\
Moves-left & $M = 86.6$ plies & expects the game to last another 43 moves \\
\bottomrule
\end{tabular}
\end{center}
All three heads are consistent with each other --- and all three are wrong about a position
that ends in two moves. The search then found it, in 188 nodes, by the mechanism of
Chapter~\ref{ch:puct} §6.5.
\end{realdata}

\begin{keyidea}
This position is a clean, reproducible instance of the phenomenon the suppressed-win idea is
about: \textbf{the output layer under-rates a forcing win that is objectively there}. What it
does \emph{not} yet show is anything about the middle layers. Whether layer 8 "knew" is
exactly the open question --- and it is answerable, on this position, with the machinery this
project already has.
\end{keyidea}

\section{The experiment, designed properly}

Here is the suppressed-win experiment written out so that it could fail.

\subsection{The quantity}

Train a probe (Chapter~\ref{ch:probes}) at layer $l$ to predict a property such as
\emph{"there exists a forced win for the side to move within $k$ plies"}. Write its output
$\hat{y}^{(l)}(s) \in [0,1]$. Define
\begin{equation}
  \label{eq:suppression}
  \Delta_{\text{suppress}}(s) \;=\; \hat{y}^{(l)}(s) \;-\; \max_{a \in \mathcal{W}(s)} P(a\given s),
\end{equation}
where $\mathcal{W}(s)$ is the set of moves that actually achieve the win. Large positive
$\Delta$ means: the middle layer indicates a win is available while the policy assigns little
probability to any move that takes it.

\subsection{The controls, without which the number is noise}

\begin{enumerate}
\item \textbf{Input baseline.} The identical probe on layer 0. If a linear function of the
  raw board detects "a forced win exists" nearly as well, the middle layers have added nothing.
\item \textbf{Untrained-network baseline.} The identical probe on a randomly initialised net.
  This catches probes that are exploiting the geometry of the encoding rather than learned
  content.
\item \textbf{Label-shuffle control.} Retrain on shuffled labels; accuracy must collapse to the
  base rate. If it does not, there is leakage --- most likely positions from the same game
  appearing in both splits.
\item \textbf{Causal test.} Ablate or patch the implicated activations and show the policy
  changes. Without this you have shown correlation, and the whole claim is about a
  \emph{disagreement} between two parts of the network, which correlation cannot establish.
\item \textbf{Held-out re-test.} Whatever layer wins across the 15 candidates, re-measure it on
  positions never used in the search (Chapter~\ref{ch:probes}, Trap 4).
\end{enumerate}

\subsection{The decisive test}

Even with all five controls, one confound remains, and it is the one that would fool a careful
person. A probe trained to detect "a forced win exists" may simply be detecting
\textbf{tactical density} --- loose pieces, exposed king, forcing checks available --- which
correlates with forced wins and is surely represented in a strong network. Such a probe would
fire on many sharp positions with no win at all.

\begin{keyidea}
\textbf{The decisive test is matched pairs.} Construct pairs of positions that are as similar
as possible --- same material, same structure, same tactical density --- in one of which the
forcing win exists and in the other of which it does not (one defensive resource added, one
tempo different). A probe reading a genuine look-ahead representation separates the pair. A
probe reading "this looks tactical" does not. Report accuracy \emph{within pairs}; it is the
only number that means what you want it to mean.
\end{keyidea}

\section{Base rates: why the naive detector fails as a feature}

Suppose all of the above succeeds and you have a probe with genuinely good numbers. There is
still an arithmetic obstacle between a working probe and a usable feature, and it is worth
doing the sums before building anything.

\begin{worked}{the false-alarm calculation}
Suppose you scan positions where the policy's top move is not a decisive tactic, looking for
suppressed wins. Assume, generously, that \textbf{1 in 500} such positions really does hide a
low-policy forced win. Suppose the probe is good: it catches $90\%$ of the real cases (recall
$0.9$) and fires on only $5\%$ of the rest (false-positive rate $0.05$).

Out of 10{,}000 scanned positions:
\begin{itemize}
\item real suppressed wins: $10{,}000/500 = 20$. The probe catches $0.9 \times 20 = 18$.
\item non-cases: $9{,}980$. The probe fires on $0.05 \times 9{,}980 = 499$.
\end{itemize}
So of $18 + 499 = 517$ alerts, $18$ are real:
\[
  \text{precision} \;=\; \frac{18}{517} \;=\; 3.5\% .
\]
\textbf{Thirty-five out of every thirty-six alerts would be wrong.} And each wrong alert is
the tool telling you, with the full authority of "the engine's internals", that a brilliancy
was available when it was not.
\end{worked}

\begin{keyidea}
This is the base-rate problem, and it is not solved by improving the probe a little. To reach
$50\%$ precision at a $1/500$ base rate you need a false-positive rate near $0.2\%$ --- a
$25\times$ improvement, on a rare and subtle property. \textbf{The binding constraint on this
feature is precision, not accuracy}, and any evaluation reported as "accuracy" is measuring
the wrong thing.
\end{keyidea}

\subsection{What actually makes it feasible}

The good news is that you do not have to run the probe on raw positions. You already have a
cheap, reliable filter --- \emph{the search itself}:

\begin{enumerate}
\item Run a search. Identify positions where a move with low policy turns out best, or has a
  proven value (Chapter~\ref{ch:readingtree}'s corrected gem criterion). This filter is
  \textbf{verified}: the win is established by search and confirmable by Stockfish, so there
  are no false positives about the chess.
\item \emph{Then} ask the probe its question --- not "is there a win?" but "did the middle
  layers represent it?"
\end{enumerate}

This reordering changes the product. The user-facing claim becomes:

\begin{center}
\begin{tabular}{@{}p{6.3cm}p{6.9cm}@{}}
\toprule
\textbf{Naive design} & \textbf{Filtered design} \\
\midrule
"The network's deeper layers saw a winning sacrifice here." & "Search found a winning move
here that the network's instinct rated at $1.6\%$." \\
\addlinespace
Rests on the probe. Precision $3.5\%$. Unfalsifiable to the user. & Rests on the search.
Precision $\approx 100\%$. Checkable. \\
\addlinespace
The probe is load-bearing. & The probe is a \emph{bonus} claim, added only when it is
confident, and removable without breaking the feature. \\
\bottomrule
\end{tabular}
\end{center}

\begin{keyidea}
\textbf{Put the verified thing in the load-bearing position and the speculative thing on
top.} The training value to you --- here is a strong move your instinct would not have
considered, now play it out --- comes entirely from the search-verified part. The
interpretability claim adds a mechanism story, which is fascinating and should be presented as
what it is. If the probe turns out not to work on BT3, the feature still ships.
\end{keyidea}

\section{What this means for the Tal work}

Your project's interest here is not idle: the sacrificial, sharp chess you want to play
\emph{is} the region where a policy trained on outcomes is most likely to be conservative, and
the region where you personally report a specific fear --- not being able to foresee the
position a sacrifice produces.

Two things follow from this chapter.

\begin{enumerate}
\item \textbf{The search-verified version of this feature is available now.} Nothing in it
  requires probes. Chapter~\ref{ch:readingtree} gives the detector, corrected for the
  proven-value case; the Morphy position is a working example with real numbers; the drill is
  to play the continuation out.
\item \textbf{The interpretability version is a research project with a specific first
  experiment.} Probe layer 8 of BT3 on matched pairs, with all five controls, and see whether
  the pairs separate. On present evidence it is a good bet and an unproven one.
\end{enumerate}

\begin{hypothesis}
To be explicit about the status of the central idea: \emph{that BT3's middle layers represent
a winning tactic which its policy head then suppresses} is, in this project, an untested
hypothesis. It is motivated by published work on a related network, and by a real observed
failure of the output heads (§15.2). It has not been measured here. Everything above is the
design of the measurement.
\end{hypothesis}

\begin{repobox}[title={in this repository}]
The pieces exist: \texttt{backend/neural\_vision.py} for hooks, the corpus for positions,
Stockfish for ground truth, and the collection scripts in this book's \texttt{tools/} for
harvesting policy-versus-search divergences at scale. The missing piece is the labelled
dataset of matched pairs --- which is a chess-generation problem, not a machine-learning one,
and is the natural next task if this line is pursued.
\end{repobox}

\begin{exercises}

\exhead[stating the claim]{ch15:claim} Write down, in one sentence each, the three findings of
the published look-ahead work, and mark which of them requires a causal intervention to
support.

\exhead[architecture specificity]{ch15:arch} Give two concrete reasons a result obtained on
one Leela network might not hold on BT3. What is the cheapest experiment that would check?

\exhead[the position]{ch15:position} Using the real numbers in §15.2, write the three
statements that are safe to make about this position and the one tempting statement that is
not yet supported.

\exhead[matched pairs]{ch15:pairs} Design a matched pair by hand: describe two positions
differing minimally, one with a forced win and one without. What is the smallest difference
you can make that removes the win? Why does the size of that difference matter for the test?

\exhead[base rates]{ch15:baserate} Redo the calculation of §15.4 with a base rate of $1/100$
and a false-positive rate of $0.02$. What precision do you get? Now find the false-positive
rate needed for $80\%$ precision at a $1/500$ base rate.

\exhead[the wrong metric]{ch15:wrongmetric} A report says the suppressed-win probe is "$94\%$
accurate". At a base rate of $1/500$, what accuracy does a probe achieve by always answering
"no"? What should have been reported instead?

\exhead[reordering]{ch15:reorder} Explain in your own words why filtering with search first
and probing second changes the feature's precision so dramatically. Which of the two designs
would survive the probe being subtly broken, and why does that matter?

\exhead[the user-facing sentence]{ch15:sentence} Write the exact text your tool should show
for the Morphy position, in the filtered design. Then write the version that would be
justified \emph{if} the probe experiment succeeded, marking clearly what extra claim it makes.

\exhead[falsification]{ch15:falsify} State the result that would make you abandon the
suppressed-win probe. Be specific: which number, below which value, on which test set. (An
idea you cannot specify a kill condition for is not yet an experiment.)

\exhead[the citation]{ch15:citation} §15.1 flags a citation this book declined to rely on.
Describe how you would verify it, and write the note you would add to the project's research
document in each of the two possible outcomes.

\end{exercises}
